\documentclass[11pt]{article}

\usepackage{amsmath,amssymb,amsthm,mathtools}
\usepackage[margin=1in]{geometry}
\usepackage{longtable}
\usepackage{url}

\newtheorem{theorem}{Theorem}[section]
\newtheorem{lemma}[theorem]{Lemma}
\newtheorem{corollary}[theorem]{Corollary}

\begin{document}

\section{The singular hub family}

For a finite graph $H$, write
\[
 \Gamma_H(W)=t(H,W)-\Phi_H(p(W)),
 \qquad
 \Phi_H(p)=(1-p)^{v(H)}
 \chi_H\!\left(\frac1{1-p}\right),
\]
with the second expression interpreted by polynomial continuation at $p=1$.
This note treats a full three-parameter singular family in the requested
high-density test.  Split an arbitrary probability space into measurable
sets $A,B$ of masses $s,1-s$.  For $a,r\in[0,1]$, define
\begin{equation}
 W_{s,a,r}(x,y)=
 \begin{cases}
  1-r,&x,y\in A,\\
  1-a,&x\in A, y\in B\text{ or conversely},\\
  1,&x,y\in B.
 \end{cases}
 \label{eq:hub-family}
\end{equation}
Thus the complement is supported on pairs incident to the exceptional set.
Its degree on $A$ is $(1-s)a+sr$, which may converge to the nonzero value
$a$ as $s\downarrow0$.  Therefore this family is not covered by any theorem
requiring uniformly small complement degree.  It includes:
\begin{itemize}
 \item the one-clique hole, when $(a,r)=(0,1)$;
 \item one complete clique plus dust, when $(a,r)=(1,1)$; and
 \item every constant two-block interpolation between these boundary shapes.
\end{itemize}
The result below is uniform over the complete parameter cube
$0\leq s\leq1/3$, $0\leq a,r\leq1$.


\section{Exact density polynomial}

For $S\subseteq V(H)$, let $e_H(S)$ be the number of edges with both ends in
$S$, and let $\partial_HS$ be the set of edges with exactly one end in $S$.

\begin{lemma}[Subset expansion]
\label{lem:subset}
If $n=v(H)$, then
\begin{align}
 t(H,W_{s,a,r})
 &=\sum_{S\subseteq V(H)}
 s^{|S|}(1-s)^{n-|S|}
 (1-r)^{e_H(S)}(1-a)^{|\partial_HS|},
 \label{eq:hom-density}\\
 q:=1-p(W_{s,a,r})
 &=2s(1-s)a+s^2r.
 \label{eq:complement-density}
\end{align}
If $\chi_H(x)=\sum_{j=0}^nc_jx^j$, then
\begin{equation}
 \Phi_H(1-q)=\sum_{j=0}^nc_jq^{n-j}.
 \label{eq:target-polynomial}
\end{equation}
Consequently the gap
\begin{equation}
 G_H(s,a,r)=t(H,W_{s,a,r})-\Phi_H(p(W_{s,a,r}))
 \label{eq:gap}
\end{equation}
is a polynomial in $s,a,r$ with integer coefficients, and $s$ divides it.
\end{lemma}

\begin{proof}
For a random map $f:V(H)\to\Omega$, condition on the exact vertex set
$S=f^{-1}(A)$.  The assignment probability is
$s^{|S|}(1-s)^{n-|S|}$.  Its internal-$A$ and crossing edge factors are
respectively $1-r$ and $1-a$; all $B^2$ factors equal one.  Summing the
disjoint cases proves \eqref{eq:hom-density}.  Integrating the three cells of
$1-W_{s,a,r}$ proves \eqref{eq:complement-density}.  Substitution into the
definition of $\Phi_H$ gives \eqref{eq:target-polynomial}.  At $s=0$ the
graphon is one, both homomorphism density and target equal one, and hence the
polynomial $G_H$ has the factor $s$.
\end{proof}

No graphon approximation is present in this reduction: \eqref{eq:hub-family}
is defined on the original probability space, and only the measures of its
two measurable cells enter the exact integrals.


\section{Trivariate Bernstein certificate}

Put $x=3s$ and define the polynomial
\begin{equation}
 Q_H(x,a,r)=\left.\frac{G_H(s,a,r)}s\right|_{s=x/3}.
 \label{eq:quotient}
\end{equation}
If $(d_x,d_a,d_r)$ is its multidegree, its native trivariate Bernstein
expansion on $[0,1]^3$ is
\begin{equation}
 Q_H(x,a,r)=
 \sum_{i=0}^{d_x}\sum_{j=0}^{d_a}\sum_{k=0}^{d_r}
 \beta_{i,j,k}
 \binom{d_x}{i}\binom{d_a}{j}\binom{d_r}{k}
 x^i(1-x)^{d_x-i}
 a^j(1-a)^{d_a-j}
 r^k(1-r)^{d_r-k}.
 \label{eq:bernstein}
\end{equation}
For a power-basis polynomial
$Q_H=\sum_{u,v,w}c_{u,v,w}x^ua^vr^w$, the coefficient is reconstructed by
the exact finite formula
\begin{equation}
 \beta_{i,j,k}=
 \sum_{u\leq i,\,v\leq j,\,w\leq k}
 c_{u,v,w}
 \frac{\binom iu}{\binom{d_x}u}
 \frac{\binom jv}{\binom{d_a}v}
 \frac{\binom kw}{\binom{d_r}w}.
 \label{eq:coefficient-formula}
\end{equation}
Every basis function in \eqref{eq:bernstein} is nonnegative on the cube.

\begin{theorem}[Two-block singular-hub theorem]
\label{thm:hub}
For every one of the $29$ graphs identified below and all
\[
 0\leq s\leq\frac13,
 \qquad0\leq a,r\leq1,
\]
one has
\begin{equation}
 t(H,W_{s,a,r})\geq\Phi_H(p(W_{s,a,r})).
 \label{eq:main}
\end{equation}
In particular the conclusion holds whenever the density is in the
conjectured interval.  No member of this full persistent-hub cube is a
counterexample.
\end{theorem}

\begin{proof}
Equations \eqref{eq:hom-density}--\eqref{eq:coefficient-formula} give a finite
exact reconstruction of every coefficient $\beta_{i,j,k}$.  Direct integer
and rational arithmetic gives $\beta_{i,j,k}\geq0$ for every graph and every
index.  The complete audit signatures are recorded in
Table~\ref{tab:certificates}: the number of zero coefficients plus the least
positive and final coefficients independently check the reconstructed
arrays.  Therefore $Q_H\geq0$ on $[0,1]^3$.  Since
$G_H=sQ_H(3s,a,r)$ and $s\geq0$, inequality \eqref{eq:main} follows.  At
$s=0$, or whenever $a=r=0$, the graphon is one and the polynomial endpoint
gives equality directly.
\end{proof}


\section{Graph identities, chromatic polynomials, and targets}

All graphs use vertex set $\{0,\ldots,n-1\}$.  The graph6 string and labelled
edge set make each row unambiguous.  The final column refers to the polynomial
type in the following table.

\scriptsize
\begin{longtable}{r@{\quad}lrrr@{\quad}p{8.2cm}}
Atlas&graph6&$n$&$m$&type&edge set\\ \hline
43&\verb|Dlc|&5&6&1&01,03,04,12,23,34\\
118&\verb|Eht?|&6&7&2&01,04,12,14,15,23,34\\
122&\verb|Eld?|&6&7&2&01,03,04,12,15,23,34\\
124&\verb|Exd?|&6&7&2&01,02,04,12,15,23,34\\
126&\verb|ERUO|&6&7&2&02,05,13,14,23,34,35\\
127&\verb|EZEG|&6&7&3&02,05,12,13,23,34,45\\
130&\verb|E{CW|&6&7&4&01,02,03,12,34,35,45\\
137&\verb|EjsG|&6&8&5&01,04,12,13,14,23,34,45\\
139&\verb|Eju?|&6&8&5&01,04,05,12,13,14,23,34\\
145&\verb|Exe_|&6&8&6&01,02,04,05,12,23,25,34\\
147&\verb|EhMg|&6&8&6&01,05,12,23,24,25,34,45\\
148&\verb|EyUG|&6&8&6&01,02,05,12,13,14,34,45\\
151&\verb|EhdW|&6&8&7&01,04,12,15,23,34,35,45\\
152&\verb|EhNG|&6&8&6&01,05,12,15,23,24,34,45\\
153&\verb|Ehf_|&6&8&8&01,04,05,12,15,23,25,34\\
157&\verb|E~`_|&6&9&9&01,02,03,04,12,13,15,23,25\\
160&\verb|E~@g|&6&9&9&01,02,03,12,13,15,23,25,45\\
168&\verb|E^MG|&6&9&10&02,03,05,12,13,23,24,34,45\\
169&\verb|Exf_|&6&9&11&01,02,04,05,12,15,23,25,34\\
171&\verb|Ehew|&6&9&12&01,04,05,12,23,25,34,35,45\\
174&\verb|EtTg|&6&9&13&01,02,03,14,15,23,25,34,45\\
178&\verb|En}?|&6&10&14&01,03,04,05,12,13,14,23,24,34\\
181&\verb|E^mG|&6&10&15&02,03,04,05,12,13,23,24,34,45\\
185&\verb|Exv_|&6&10&16&01,02,04,05,12,14,15,23,25,34\\
188&\verb|EzNG|&6&10&17&01,02,05,12,13,15,23,24,34,45\\
194&\verb|E~wW|&6&11&18&01,02,03,04,12,13,14,23,24,35,45\\
196&\verb|ER~g|&6&11&18&02,04,05,13,14,15,23,24,25,34,45\\
199&\verb|El^g|&6&11&19&01,03,05,12,14,15,23,24,25,34,45\\
203&\verb|E~^G|&6&12&20&01,02,03,05,12,13,14,15,23,24,34,45
\end{longtable}
\normalsize

Types $1$--$8$, $10$, $12$, $13$, and $17$ have chromatic number three;
the other types have chromatic number four.  Deletion--contraction gives:

\scriptsize
\begin{longtable}{r@{\quad}p{7.1cm}p{7.1cm}}
type&$\chi_H(x)$&$\Phi_H(p)$\\ \hline
1&$x(x-1)(x-2)(x^2-3x+3)$&$p(2p-1)(3p^2-3p+1)$\\
2&$x(x-1)^2(x-2)(x^2-3x+3)$&$p^2(2p-1)(3p^2-3p+1)$\\
3&$x(x-1)(x-2)^2(x^2-2x+2)$&$p(2p-1)^2(2p^2-2p+1)$\\
4&$x(x-1)^3(x-2)^2$&$p^3(2p-1)^2$\\
5&$x(x-1)^2(x-2)^3$&$p^2(2p-1)^3$\\
6&$x(x-1)(x-2)^2(x^2-3x+3)$&$p(2p-1)^2(3p^2-3p+1)$\\
7&$x(x-1)(x-2)^2(x^2-3x+4)$&$p(2p-1)^2(4p^2-5p+2)$\\
8&$x(x-1)(x-2)(x^3-5x^2+9x-7)$&$p(2p-1)(7p^3-12p^2+8p-2)$\\
9&$x(x-1)^2(x-2)^2(x-3)$&$p^2(2p-1)^2(3p-2)$\\
10&$x(x-1)(x-2)^2(x^2-4x+5)$&$p(2p-1)^2(5p^2-6p+2)$\\
11&$x(x-1)(x-2)(x-3)(x^2-3x+3)$&$p(2p-1)(3p-2)(3p^2-3p+1)$\\
12&$x(x-1)(x-2)(x^3-6x^2+13x-11)$&$p(2p-1)(11p^3-20p^2+13p-3)$\\
13&$x(x-1)(x-2)(x^3-6x^2+14x-13)$&$p(2p-1)(13p^3-25p^2+17p-4)$\\
14&$x(x-1)^2(x-2)(x-3)^2$&$p^2(2p-1)(3p-2)^2$\\
15&$x(x-1)(x-2)^3(x-3)$&$p(2p-1)^3(3p-2)$\\
16&$x(x-1)(x-2)(x-3)(x^2-4x+5)$&$p(2p-1)(3p-2)(5p^2-6p+2)$\\
17&$x(x-1)(x-2)(x^3-7x^2+18x-17)$&$p(2p-1)(17p^3-33p^2+22p-5)$\\
18&$x(x-1)(x-2)(x-3)(x^2-5x+7)$&$p(2p-1)(3p-2)(7p^2-9p+3)$\\
19&$x(x-1)(x-2)(x-3)(x^2-5x+8)$&$p(2p-1)(3p-2)(8p^2-11p+4)$\\
20&$x(x-1)(x-2)(x-3)(x^2-6x+10)$&$p(2p-1)(3p-2)(10p^2-14p+5)$
\end{longtable}
\normalsize


\section{Exact certificate audit}

Table~\ref{tab:certificates} gives the multidegree, number $z_H$ of zero
Bernstein coefficients, least positive coefficient $m_H$, and final
coefficient $f_H=\beta_{d_x,d_a,d_r}$.  Every unreported coefficient is
nonnegative; the verifier reconstructs the complete array from
\eqref{eq:coefficient-formula}.  The arrays contain between $336$ and $1300$
entries, $22{,}246$ in total, so these four independent signatures are more
reviewable than printing every rational number.

\begin{table}[ht]
\centering
\scriptsize
\begin{tabular}{r|c|r|c|c@{\qquad}r|c|r|c|c}
$H$&$(d_x,d_a,d_r)$&$z_H$&$m_H$&$f_H$&
$H$&$(d_x,d_a,d_r)$&$z_H$&$m_H$&$f_H$\\ \hline
43&(7,5,6)&40&$1/315$&$316/729$&118&(9,6,7)&49&$1/567$&$1840/6561$\\
122&(9,6,7)&49&$11/6804$&$1840/6561$&124&(9,6,7)&49&$11/6804$&$1840/6561$\\
126&(9,6,7)&49&$11/6804$&$1840/6561$&127&(9,6,7)&49&$5/3402$&$5020/19683$\\
130&(9,5,7)&49&$5/3402$&$5120/19683$&137&(9,6,8)&52&$17/9072$&$5200/19683$\\
139&(9,6,8)&52&$1/567$&$5200/19683$&145&(9,6,8)&52&$1/567$&$1700/6561$\\
147&(9,7,8)&52&$5/3024$&$1700/6561$&148&(9,6,8)&52&$5/3024$&$1700/6561$\\
151&(9,7,8)&52&$1/648$&$5000/19683$&152&(9,6,8)&52&$1/648$&$1700/6561$\\
153&(9,6,8)&52&$1/648$&$4600/19683$&157&(9,7,9)&55&$11/5832$&$1760/6561$\\
160&(9,7,9)&55&$7/3888$&$1760/6561$&168&(9,8,9)&55&$5/2916$&$5080/19683$\\
169&(9,7,9)&55&$5/2916$&$520/2187$&171&(9,7,9)&55&$19/11664$&$4580/19683$\\
174&(9,7,9)&55&$1/648$&$4480/19683$&178&(9,7,10)&58&$1/540$&$640/2187$\\
181&(9,8,10)&58&$13/7290$&$1720/6561$&185&(9,8,10)&58&$5/2916$&$1520/6561$\\
188&(9,8,10)&58&$2/1215$&$4460/19683$&194&(9,8,11)&61&$31/17820$&$1480/6561$\\
196&(9,8,11)&61&$31/17820$&$1480/6561$&199&(9,9,11)&61&$1/594$&$1280/6561$\\
203&(9,9,12)&64&$37/21384$&$1240/6561$&&&&&
\end{tabular}
\caption{Exact native trivariate Bernstein signatures.}
\label{tab:certificates}
\end{table}

Run
\begin{verbatim}
python codes/principled_negative_tests.py
\end{verbatim}
from the repository root.  The line headed ``singular two-block hub
full-cube theorem'' reconstructs every graph from the Atlas, enumerates every
vertex subset, recomputes its chromatic polynomial, performs the exact
power-to-Bernstein conversion, checks every coefficient sign, and compares
the displayed signatures with hard-coded rational values.


\section{Endpoint, assumptions, and exact boundary}

At $p=1$, formula \eqref{eq:complement-density} forces either $s=0$ or
$a=r=0$ (apart from irrelevant zero-mass cells).  Hence $W=1$ almost
everywhere, $t(H,W)=1$, and monicity of $\chi_H$ gives
$\Phi_H(1)=1$.  No division by $1-p$ is used in the proof.

Every calculation uses bounded measurable cell indicators and finite sums,
so Fubini is automatic.  The underlying probability space need not be
atomless.  There is no regularity, finite-graph approximation, optimization,
or tensor construction.  The family itself has two measurable blocks, but
the theorem is exhaustive over all three of its continuous parameters; no
direction inside that family has been selected after inspection.

The exact boundary is important.  This theorem assumes the complement values
are constant on the cut and on $A^2$, and does not cover exceptional mass
above $1/3$.  The independent theorem in
\path{notes/arbitrary_fibre_hub_high_density_stability.tex} first allows
completely arbitrary cut fibres and exceptional-square kernels under an
adaptive $L^2$-energy condition and then removes that condition on a smaller
uniform interval.  It covers any number of interacting hubs inside the same
small exceptional set.  The further high-degree-set theorem in
\path{notes/arbitrary_graphon_high_density_stability.tex} combines that
margin with maximum-degree stability and closes the endpoint for every
graphon.  Intermediate densities remain unresolved, so no catalogue status
changes.

For formal verification, the analytic layer needs only
Lemma~\ref{lem:subset} and Bernstein nonnegativity.  The finite layer should
certify the vertex-subset signatures, chromatic polynomials, and complete
coefficient arrays for the $29$ rows.  Since this is a restricted family, no
Atlas status module should be changed.

\end{document}
